\clearpage
\section{Результаты работы}

Для проверки прямого и обратного преобразования использовано число $11{,}55$.
Программа сформировала представление

\[
    11{,}55_{10}\longrightarrow
    \underbrace{0}_{\text{знак}}\;
    \underbrace{10000010}_{\text{экспонента}}\;
    \underbrace{0111001}_{\text{мантисса}}.
\]

Результат прямого и обратного преобразований приведён в
листинге~\ref{lst:conversion-example}.

\codeListing[text][fontsize=\footnotesize]
    {listings/conversion-example.txt}
    {Преобразование числа $11{,}55$ в формат bfloat16 и обратно}
    {\label{lst:conversion-example}}

При декодировании получено значение $11{,}5625$. Абсолютная ошибка равна
$0{,}0125$ и обусловлена округлением семибитной мантиссы.

\newpage
